\documentclass[11pt, a4paper]{article}
\usepackage[utf8]{inputenc}
\usepackage{amsmath, amssymb}
\usepackage{graphicx}
\usepackage{hyperref}
\usepackage{xcolor}
\usepackage{listings}
\lstset{
    language=Python,
    basicstyle=\ttfamily\small,
    keywordstyle=\color{blue}\bfseries,
    commentstyle=\color{gray}\itshape,
    stringstyle=\color{orange},
    numberstyle=\tiny\color{gray},
    numbers=left,
    stepnumber=1,
    numbersep=8pt,
    frame=single,
    breaklines=true,
    breakatwhitespace=false,
    showstringspaces=false,
    tabsize=4,
    captionpos=b
}
\usepackage{geometry}
\geometry{margin=1in}

\title{EE 325 Digital Signal Processing \\ Assignment 2: Time Domain Analysis of LTI Systems}
\author{PERERA J.D.T. \\ E/21/291}
\date{\today}

\begin{document}
\maketitle

\section{Objective}
Analyze the time-domain response of a Linear Time-Invariant (LTI) system using difference equations in Python.

\section{System Description}
The system is defined by the first-order difference equation:
$$y[n] = 0.5 y[n-1] + x[n]$$
Rearranging terms, this corresponds to the system:
$$y[n] - 0.5 y[n-1] = x[n]$$
Impulse Response expectation: $h[n] = (0.5)^n u[n]$.

\section{Simulation Results}
\subsection{1. Impulse Response ($h[n]$)}
\begin{itemize}
    \item \textbf{Input:} $x[n] = \delta[n]$ (Unit Impulse).
    \item \textbf{Method:} Iteratively inferred $y[n]$ with $y[-1]=0$.
    \item \textbf{Result:}
    \begin{itemize}
        \item $n=0: y[0] = 0.5(0) + 1 = 1$
        \item $n=1: y[1] = 0.5(1) + 0 = 0.5$
        \item $n=2: y[2] = 0.5(0.5) + 0 = 0.25$
    \end{itemize}
    \item \textbf{Verification:} Follows the theoretical decay $h[n] = (0.5)^n$.
\end{itemize}

\subsection{2. Step Response ($s[n]$)}
\begin{itemize}
    \item \textbf{Input:} $x[n] = u[n]$ (Unit Step).
    \item \textbf{Result:} Accumulates the impulse response.
    \begin{itemize}
        \item $n=0: 1$
        \item $n=1: 1 + 0.5 = 1.5$
        \item $n=2: 1.5 + 0.25 = 1.75$
    \end{itemize}
    \item \textbf{Steady State:} As $n \to \infty$, $s[n] \to \frac{1}{1 - 0.5} = 2$.
\end{itemize}

\section{Conclusion}
The simulation successfully implements the recursive difference equation and visualizes the system's characteristic exponential decay (Stable IIR system).

\appendix
\section{Python Source Code}
\textbf{File:} \texttt{simulation.py}
\begin{lstlisting}[caption={Time Domain LTI System Simulation -- Assignment 2}]
import numpy as np
import matplotlib.pyplot as plt

# Discrete time indices
n = np.arange(0, 21)

# Define input signals
delta_n = np.zeros(21)  
delta_n[0] = 1

u_n = np.ones(21)  

#  y[n] = 0.5*y[n-1] + x[n]
a1 = 0.5
b0 = 1.0  

# Initialize output arrays
h_n = np.zeros(21)  
s_n = np.zeros(21)  

# Simulate impulse response h[n]
for i in range(21):
    if i == 0:
        h_n[i] = b0 * delta_n[i]
    else:
        h_n[i] = a1 * h_n[i-1] + b0 * delta_n[i]

# Simulate step response s[n]
for i in range(21):
    if i == 0:
        s_n[i] = b0 * u_n[i]
    else:
        s_n[i] = a1 * s_n[i-1] + b0 * u_n[i]


# Plotting the results

#IMPULSE RESPONSE PLOT
plt.figure(figsize=(12, 8))
plt.subplot(2, 1, 1)
plt.stem(n, h_n, linefmt='b-', markerfmt='bo', basefmt=' ', label='Impulse Response h[n]')
plt.title('Impulse Response of the System') 
plt.xlabel('n (Discrete Time Index)')
plt.ylabel('h[n]')
plt.legend()
plt.grid()

#STEP RESPONSE PLOT
plt.subplot(2, 1, 2)
plt.stem(n, s_n, linefmt='r-', markerfmt='ro', basefmt=' ', label='Step Response s[n]')
plt.title('Step Response of the System')
plt.xlabel('n (Discrete Time Index)')
plt.ylabel('s[n]')
plt.legend()
plt.grid()
plt.tight_layout()
plt.show()
\end{lstlisting}

\end{document}
